\documentclass[10pt,a4paper]{amsart}
\usepackage{amsmath,amssymb,amsthm,graphicx}
% Retained minimal preamble for the Stacks Project chapter more-algebra.tex (More on Algebra).
% The chapter file pinned-original-source.tex carries only \input{preamble} on its first line, so no
% preamble text lives inside the chapter itself.  The Stacks Project uses one single shared preamble for
% every chapter; the snapshot stages it identically as sources/stacks-chow/preamble.tex and
% sources/stacks-etale-cohomology/preamble.tex.  The block below is copied mechanically, line for line,
% from sources/stacks-chow/preamble.tex lines 14-18 (the Xy-pic setup) and 154-217 (the native theorem
% declarations together with the mathematical macros).  Nothing was added, removed or redefined; the
% wrapper class and the amsmath/amssymb/amsthm/graphicx base packages are supplied by the assembler.
% Chapter: more-algebra.tex -- More on Algebra.
% For commutative diagrams we use Xy-pic
\usepackage[all]{xy}

% We use 2cell for 2-commutative diagrams.
\xyoption{2cell}

% Theorem environments.
%
\theoremstyle{plain}
\newtheorem{theorem}[subsection]{Theorem}
\newtheorem{proposition}[subsection]{Proposition}
\newtheorem{lemma}[subsection]{Lemma}

\theoremstyle{definition}
\newtheorem{definition}[subsection]{Definition}
\newtheorem{example}[subsection]{Example}
\newtheorem{exercise}[subsection]{Exercise}
\newtheorem{situation}[subsection]{Situation}

\theoremstyle{remark}
\newtheorem{remark}[subsection]{Remark}
\newtheorem{remarks}[subsection]{Remarks}

\numberwithin{equation}{subsection}

% Macros
%
\def\lim{\mathop{\mathrm{lim}}\nolimits}
\def\colim{\mathop{\mathrm{colim}}\nolimits}
\def\Spec{\mathop{\mathrm{Spec}}}
\def\Hom{\mathop{\mathrm{Hom}}\nolimits}
\def\Ext{\mathop{\mathrm{Ext}}\nolimits}
\def\SheafHom{\mathop{\mathcal{H}\!\mathit{om}}\nolimits}
\def\SheafExt{\mathop{\mathcal{E}\!\mathit{xt}}\nolimits}
\def\Sch{\mathit{Sch}}
\def\Mor{\mathop{\mathrm{Mor}}\nolimits}
\def\Ob{\mathop{\mathrm{Ob}}\nolimits}
\def\Sh{\mathop{\mathit{Sh}}\nolimits}
\def\NL{\mathop{N\!L}\nolimits}
\def\CH{\mathop{\mathrm{CH}}\nolimits}
\def\proetale{{pro\text{-}\acute{e}tale}}
\def\etale{{\acute{e}tale}}
\def\QCoh{\mathit{QCoh}}
\def\Ker{\mathop{\mathrm{Ker}}}
\def\Im{\mathop{\mathrm{Im}}}
\def\Coker{\mathop{\mathrm{Coker}}}
\def\Coim{\mathop{\mathrm{Coim}}}

% Boxtimes
%
\DeclareMathSymbol{\boxtimes}{\mathbin}{AMSa}{"02}

%
% Macros for moduli stacks/spaces
%
\def\QCohstack{\mathcal{QC}\!\mathit{oh}}
\def\Cohstack{\mathcal{C}\!\mathit{oh}}
\def\Spacesstack{\mathcal{S}\!\mathit{paces}}
\def\Quotfunctor{\mathrm{Quot}}
\def\Hilbfunctor{\mathrm{Hilb}}
\def\Curvesstack{\mathcal{C}\!\mathit{urves}}
\def\Polarizedstack{\mathcal{P}\!\mathit{olarized}}
\def\Complexesstack{\mathcal{C}\!\mathit{omplexes}}
% \Pic is the operator that assigns to X its picard group, usage \Pic(X)
% \Picardstack_{X/B} denotes the Picard stack of X over B
% \Picardfunctor_{X/B} denotes the Picard functor of X over B
\def\Pic{\mathop{\mathrm{Pic}}\nolimits}
\def\Picardstack{\mathcal{P}\!\mathit{ic}}
\def\Picardfunctor{\mathrm{Pic}}
\def\Deformationcategory{\mathcal{D}\!\mathit{ef}}
\usepackage{cleveref}
\crefname{originallemma}{original lemma}{original lemmas}
\Crefname{originallemma}{Original lemma}{Original lemmas}
\makeatletter
\newlabel{more-algebra-lemma-resolution-filtered-complex}{{0H7X}{source}{}{sourcepointer-1}{}}
\newlabel{more-algebra-lemma-resolution-filtered-complex@cref}{{[originallemma][0][]0H7X}{[1][0][]source}{}{}{}}
\makeatother
\crefname{originallemma}{original lemma}{original lemmas}
\Crefname{originallemma}{Original lemma}{Original lemmas}
\makeatletter
\newlabel{lemma-colimit-K-flat}{{06Y3}{source}{}{sourcepointer-2}{}}
\newlabel{lemma-colimit-K-flat@cref}{{[originallemma][0][]06Y3}{[1][0][]source}{}{}{}}
\makeatother
\crefname{originallemma}{original lemma}{original lemmas}
\Crefname{originallemma}{Original lemma}{Original lemmas}
\makeatletter
\newlabel{lemma-K-flat-two-out-of-three-ses}{{0BYH}{source}{}{sourcepointer-3}{}}
\newlabel{lemma-K-flat-two-out-of-three-ses@cref}{{[originallemma][0][]0BYH}{[1][0][]source}{}{}{}}
\makeatother
\title{A bounded filtered complex is resolved by filtered complexes in the prescribed range}
\author{The Stacks Project Authors}
\date{stacks-project:more-algebra.tex:local-snapshot-20261012; GNU FDL 1.2 or later}
\begin{document}
\maketitle
\setcounter{subsection}{1}
\noindent\textit{Editorial note.} All delivered mathematical text is the author's own, copied from the pinned Stacks Project chapter file more-algebra.tex; no statement or proof step was rewritten, reordered or removed.\par\medskip
\noindent\textit{Editorial note.} The chapter file pinned-original-source.tex carries only \textbackslash{}input\{preamble\} on its first line, so no preamble text lives inside the chapter. The retained file sources/stacks-more-algebra/retained-preamble.tex is a mechanical, line-for-line copy of the project's single shared preamble as staged in the snapshot (Xy-pic setup plus theorem declarations and mathematical macros). Nothing was added, removed or redefined, and the excerpt serves only as preamble, never as a span coordinate.\par\medskip
\noindent\textit{Editorial note.} Both counted spans start at the \textbackslash{}begin line of their environment, so the statement and the author's proof are delivered whole; surrounding prose of the section is deliberately outside the counted spans and is carried as standard prerequisites instead.\par\medskip
\noindent\textit{Editorial note.} The citation fields (source\_id, source\_version, authors, source\_url, license.name, source\_title) are identical to those of every other card this lane has produced from this chapter, so the batch selector treats them as one source.\par\medskip
\noindent\textit{Allowed standard prerequisites (editorial).}\begin{itemize}
\item Filtered complex and its associated graded complexes. defined in the same section, source lines 16330-16420, and used throughout the proof. Supplies the object of the statement and the shape of its filtration.
\item A module admits a resolution by free or projective modules, and such a resolution lifts along a filtration. used at proof lines 15670-15720 to produce the next stage of the induction. Provides the complex that is glued on at each step.
\item The long exact sequence of cohomology attached to a short exact sequence of complexes. used to compare the cohomology of consecutive stages at proof lines 15700-15760. Shows that each stage is quasi-isomorphic to the original complex.
\end{itemize}\medskip
\ifdefined\hypertarget\hypertarget{sourcepointer-1}{}\fi
\noindent\textit{Source reference.} Original lemma 0H7X: focused-seeds/AGP/sources/stacks-snapshot-20261011/tags.txt, source lines 21472--21472. The lemma's own canonical tag, quoted so that readers can resolve the label inside the delivered text; it points at the counted statement itself and carries no proof content.\par\medskip
\ifdefined\hypertarget\hypertarget{sourcepointer-2}{}\fi
\noindent\textit{Source reference.} Original lemma 06Y3: focused-seeds/AGP/sources/stacks-snapshot-20261011/tags.txt, source lines 8872--8872. Navigation pointer to a statement the counted proof refers to but does not contain; nothing load-bearing is delegated to it.\par\medskip
\ifdefined\hypertarget\hypertarget{sourcepointer-3}{}\fi
\noindent\textit{Source reference.} Original lemma 0BYH: focused-seeds/AGP/sources/stacks-snapshot-20261011/tags.txt, source lines 15015--15015. Navigation pointer to a statement the counted proof refers to but does not contain; nothing load-bearing is delegated to it.\par\medskip
\begin{lemma}
\label{lemma-resolution-filtered-complex}
Let $R$ be a ring. Let $K^\bullet$ be a filtered complex. There exists a
map $f : P^\bullet \to K^\bullet$ of filtered complexes such that
\begin{enumerate}
\item each $P^n$, $F^iP^n$, $\text{gr}^iP^n$ is a free $R$-module,
\item the complexes of $R$-modules $P^\bullet$, $F^iP^\bullet$, and
$\text{gr}^iP^\bullet$ are K-flat,
\item $f$ induces quasi-isomorphisms
$P^\bullet \to K^\bullet$,
$F^iP^\bullet \to F^iK^\bullet$, and
$\text{gr}^iP^\bullet \to \text{gr}^iK^\bullet$.
\end{enumerate}
\end{lemma}

\begin{proof}
Let us say a filtered complex $L^\bullet$ is {\it basic} if
each $L^n$, $F^iL^n$, $\text{gr}^iL^n$ is a free $R$-module
and if all differentials are zero.

\medskip\noindent
There exists a basic filtered complex $P_0^\bullet$ and a map
$f_0 : P_0^\bullet \to K^\bullet$ of filtered complexes such that
$f_0$ and $F^if_0$, $i \in \mathbf{Z}$ are surjective on cohomology.
To see this set
$$
P_0^n = \bigoplus\nolimits_{z \in \Ker(d_K^n)} R \xi_z
\oplus
\bigoplus\nolimits_{j \in \mathbf{Z}}
\bigoplus\nolimits_{z \in \Ker(d_{F^jK}^n)} R \xi_z
$$
with zero differentials and the map $f_0$ defined by
$f_0(\xi_z) = z$. As for the filtration, we set
$$
F^iP_0^n = \bigoplus\nolimits_{j \geq i \in \mathbf{Z}}
\bigoplus\nolimits_{z \in \Ker(d_{F^jK}^n)} R \xi_z
$$
We leave it to the reader to check that this gives
$f_0 : P_0^\bullet \to K^\bullet$ as claimed.

\medskip\noindent
By induction on $m \geq 0$ we are going to construct embeddings
$$
P_0^\bullet \subset \ldots P_m^\bullet \subset P_{m + 1}^\bullet
$$
of filtered complexes and maps $f_m : P_m^\bullet \to K^\bullet$
with the following properties
\begin{enumerate}
\item the filtered complex $P_{m + 1}^\bullet / P_m^\bullet$ is basic,
\item the kernel of $H^n(f_m)$ and $H^n(F^if_m)$ maps to zero in
$H^n(P_{m + 1}^\bullet)$ and $H^n(F^iP_{m + 1}^\bullet)$,
\item the map $f_{m + 1} : P_{m + 1}^\bullet \to K^\bullet$
extends the map $f_m$.
\end{enumerate}
To to this, set
$$
\Omega_{n, m} = \Ker\left(\Ker(d^n_{P_m}) \to H^n(K^\bullet)\right),
\text{ resp. }
\Omega_{n, i, m} = \Ker\left(\Ker(d^n_{F^iP_m}) \to H^n(F^iK^\bullet)\right),
$$
Note that $\Omega_{n, m}$ surjects onto the kernel of $H^n(f_m)$
and that $\Omega_{n, i, m}$ surjects onto the kernel of $H^n(F^if_n)$.
For each $z \in \Omega_{n, m}$, resp.\ $z \in \Omega_{n, i, m}$
we choose a $y_z \in K^{n - 1}$, resp.\ $y_z \in F^iK^{n - 1}$
such that $d_K(y_z) = f_m(z)$. Then we set
$$
P^n_{m + 1} = P_m^n \oplus
\bigoplus\nolimits_{z \in \Omega_{n + 1, m}} R \eta_z
\oplus
\bigoplus\nolimits_{j \in \mathbf{Z}}
\bigoplus\nolimits_{z \in \Omega_{n + 1, j, m}} R \eta_z
$$
We set $d_{P_m}(\eta_z) = z$ and we set $f_{m + 1}(\eta_z) = y_z$.
Finally, we set
$$
F^iP_{m + 1}^n = F^iP_m^n \oplus
\bigoplus\nolimits_{j \geq i}
\bigoplus\nolimits_{z \in \Omega_{n + 1, j, m}} R \eta_z
$$
We leave it to the reader to check that this gives
$P_m^\bullet \subset P_{m + 1}^\bullet$ and
$f_{m + 1} : P_{m + 1}^\bullet \to K^\bullet$ as claimed.

\medskip\noindent
At this point we simply take
$$
P^\bullet = \bigcup P_m^\bullet
$$
as a filtered complex with map $f : P^\bullet \to K^\bullet$
given by $\bigcup f_m$.

\medskip\noindent
Part (1) of the statement of the lemma holds because $P^n$
as a filtered module is isomorphic to the direct sum of
$P_0^n$ and $P_{m + 1}^n/P_m^n$ for $m \geq 0$. Small detail omitted.

\medskip\noindent
Part (3) of the statement. Observe that
$H^n(P^\bullet) = \colim H^n(P_m^\bullet)$.
The map $f_0$ is surjective on cohomology; whence the same holds
for each $f_m$. For each $m$ by construction the embedding
$P_m^\bullet \subset P_{m + 1}^\bullet$
kills the kernel of $H^n(f_m)$. Combining these facts the
reader easily concludes that $H^n(f)$ is an isomorphism.
Similarly for $H^n(F^if_n)$. Then also $H^n(\text{gr}^if)$
must be an isomorphism because of the short exact sequence
$0 \to F^{i + 1} \to F^i \to \text{gr}^i \to 0$ (of functors
on the category of filtered complexes, say).
Small detail omitted.

\medskip\noindent
Part (2) of the statement. To see that $P^\bullet$ is K-flat, by
Lemma \ref{lemma-colimit-K-flat}, it suffices to show that $P_m^\bullet$
is K-flat. By Lemma \ref{lemma-K-flat-two-out-of-three-ses}
and induction it suffices to note that a complex with zero differentials
and free terms is K-flat. The same argument works to
show that $F^iP^\bullet$ is K-flat for all $i \in \mathbf{Z}$.
Finally, we see that $\text{gr}^iP^\bullet$ is K-flat by
another application of Lemma \ref{lemma-K-flat-two-out-of-three-ses}.
\end{proof}

\end{document}
